% Compilar con:  latexmk -xelatex 04-pasos-montaje.tex
\documentclass[11pt,a4paper]{article}

\usepackage{fontspec}
\usepackage[spanish,es-noshorthands]{babel}
\usepackage{geometry}
\geometry{margin=2.4cm}
\usepackage{xcolor}
\usepackage{array}
\usepackage{booktabs}
\usepackage{tabularx}
\usepackage{enumitem}
\usepackage{titlesec}
\usepackage{fancyhdr}
\usepackage[most]{tcolorbox}
\usepackage{newunicodechar}
\usepackage{hyperref}

\setmonofont{DejaVu Sans Mono}[Scale=0.82]

% --- Símbolos Unicode ---
\newunicodechar{→}{\ensuremath{\rightarrow}}
\newunicodechar{←}{\ensuremath{\leftarrow}}
\newunicodechar{↔}{\ensuremath{\leftrightarrow}}
\newunicodechar{≤}{\ensuremath{\leq}}
\newunicodechar{≥}{\ensuremath{\geq}}
\newunicodechar{×}{\ensuremath{\times}}
\newunicodechar{≈}{\ensuremath{\approx}}
\newunicodechar{·}{\textperiodcentered}
\newunicodechar{†}{\textsuperscript{\dag}}

% --- Colores y estilo ---
\definecolor{accent}{HTML}{0F6E8C}
\definecolor{accentdark}{HTML}{0A4B60}
\definecolor{soft}{HTML}{F2F6F8}
\definecolor{warn}{HTML}{7A3B00}
\definecolor{warnbg}{HTML}{FFF8F0}

\titleformat{\section}{\Large\bfseries\color{accentdark}}{\thesection}{0.6em}{}
\titleformat{\subsection}{\large\bfseries\color{accent}}{\thesubsection}{0.6em}{}
\titleformat{\subsubsection}{\normalsize\bfseries\color{accent}}{\thesubsubsection}{0.5em}{}
\titlespacing*{\section}{0pt}{1.6em}{0.5em}
\titlespacing*{\subsection}{0pt}{1.2em}{0.4em}
\titlespacing*{\subsubsection}{0pt}{0.9em}{0.3em}

\setlist[itemize]{leftmargin=1.3em,itemsep=2pt,topsep=3pt}
\setlist[enumerate]{leftmargin=1.6em,itemsep=2pt,topsep=3pt}

\renewcommand{\arraystretch}{1.25}
\setlength{\emergencystretch}{3em}
\setlength{\tabcolsep}{5pt}

\hypersetup{
  colorlinks=true, linkcolor=accent, urlcolor=accent, citecolor=accent,
  pdftitle={04 · Pasos de montaje — Runbook del sprint},
  pdfauthor={Capstone de internship · Digi International}
}

% Cajas
\newtcolorbox{nota}{colback=soft,colframe=accent!35,boxrule=0.5pt,arc=2pt,
  left=8pt,right=8pt,top=5pt,bottom=5pt}
\newtcolorbox{clave}{colback=accent!6,colframe=accent,boxrule=1pt,arc=2pt,
  left=10pt,right=10pt,top=7pt,bottom=7pt}
\newtcolorbox{aviso}{colback=warnbg,colframe=warn!60,boxrule=0.8pt,arc=2pt,
  left=8pt,right=8pt,top=5pt,bottom=5pt}

% Caja de día (título inline)
\newtcolorbox{diabox}[1]{
  colback=accent!5, colframe=accentdark, boxrule=0.8pt, arc=3pt,
  left=8pt,right=8pt,top=6pt,bottom=6pt,
  title={\bfseries\color{white}#1},
  attach boxed title to top left={yshift=-2mm,xshift=4mm},
  boxed title style={colback=accentdark,arc=2pt,boxrule=0pt}
}

% Encabezado / pie
\pagestyle{fancy}
\fancyhf{}
\fancyhead[L]{\small\color{accentdark}04 · Runbook del sprint}
\fancyhead[R]{\small\color{accentdark}Starlink-Primary Gateway Lab}
\fancyfoot[C]{\small\thepage}
\renewcommand{\headrulewidth}{0.4pt}
\renewcommand{\footrulewidth}{0pt}

\newcommand{\file}[1]{\texttt{#1}}
\newcommand{\gate}[1]{\textbf{G#1}}
\newcommand{\done}{\textbf{Done:}}
\newcommand{\embebido}{\textit{Track embebido (paralelo):}}
\newcommand{\stretchgoal}[1]{\textit{Stretch: #1}}

% ─────────────────────────────────────────────────────────────────
\begin{document}

\begin{center}
  {\Huge\bfseries\color{accentdark} Pasos de montaje}\\[6pt]
  {\large Runbook del sprint — 10 días hábiles}\\[2pt]
  {\normalsize Digi Starlink-Primary Resilient Gateway Fleet}\\[10pt]
  {\small Capstone de internship · Digi International \quad|\quad Julio 2026}
\end{center}

\vspace{4pt}
\begin{nota}
\small
Runbook \textbf{día a día} del sprint de 2 semanas (Anillo 0), con el enfoque
\textbf{Digi primero}. El track embebido corre \textbf{en paralelo} desde el Día 1 porque
los builds de Yocto son lentos y no deben bloquear el resto. Al final está la fase 2
(Opengear + expansión).

\medskip
Cada día declara: \textbf{objetivo}, \textbf{pasos}, \textbf{verificación (done)} y
\textbf{track embebido en paralelo}. Nada se marca ``done'' sin la verificación observada.
No se adivinan contratos de API: donde aparezca un gate G*, se cierra antes de depender de él
(ver \file{07-riesgos-y-verificaciones.md}).

\medskip
\textbf{Ritmo:} 10 días hábiles. Cada día trae un \textit{stretch} opcional.
Si el hardware de sitios 2-3 u Opengear llega durante el sprint, se salta a la sección de
fase 2 sin reordenar lo demás.
\end{nota}

\vspace{2pt}
{\small\tableofcontents}
\vspace{6pt}

% ─────────────────────────────────────────────────────────────────
\section{Día 0 — Charter y gates (antes de tocar hardware)}

\textbf{Objetivo:} dejar el proyecto listo para comprar/montar sin sorpresas.

\begin{enumerate}
  \item Aprobar objetivos, límites de divulgación y reglas de publicación con el manager
        (ver \file{08-executive-pitch-EN.md}).
  \item Congelar \textbf{firmware} de IX40, Hive y MP255 → matriz de versiones.
  \item \textbf{License audit (\gate{8}):} confirmar en el tenant DRM Premier + Containers +
        Push Monitor + CCCS/firmware repo.
  \item Cotizar/pedir hardware Digi del sprint y \textbf{autorizar en paralelo} sitios 2-3 y
        Opengear (fase 2).
  \item Contratar 1 plan Starlink Priority y 1 SIM 5G Carrier A; validar cobertura de banco
        (\gate{10}).
  \item Cerrar o poner en curso \gate{3} (body de \file{scan\_now}), \gate{4} (nombre del campo
        ASN remoto), \gate{9} (plan/HW Starlink), \gate{12} (\file{type}/\file{vendor\_id} del
        IX40). Los gates de Opengear (\gate{1}/\gate{2}/\gate{11}) van en paralelo, no bloquean.
\end{enumerate}

\begin{clave}
\small
\done{} matriz de versiones congelada; SKUs y licencias cotizables/reconciliados; plan Starlink
y SIM contratados; host de build DEY provisionado.
\end{clave}

% ─────────────────────────────────────────────────────────────────
\section{Semana 1 — Fundamentos Digi + arranque embebido}

% ───────────
\subsection{Día 1 — Sitio A físico + Starlink online + arranque del build Yocto}

\textbf{Objetivo:} conectividad básica del sitio y el primer \file{bitbake} corriendo.

\textbf{Pasos (in-band):}
\begin{enumerate}
  \item Rack del Sitio A, energía en PDU/UPS primaria, etiquetado.
  \item Starlink Performance Kit → IX40 \file{WAN/ETH1} (cobre). Dejar \file{SFP/ETH1} vacío.
  \item IX40 arranca; acceso local por HTTPS/SSH (no exponer a Internet).
  \item Registrar el underlay Starlink observado (\textbf{\gate{9}}): CGNAT
        \file{100.64.0.0/10} vs IPv4 pública, IPv6 \file{/64} SLAAC y \file{/56} DHCPv6-PD,
        MTU, latencia/jitter base.
\end{enumerate}

\begin{clave}
\small
\done{} el IX40 alcanza Internet por Starlink (IPv4 y, si se entrega, IPv6);
ficha del underlay registrada.
\end{clave}

\textbf{\embebido}
\begin{itemize}
  \item Provisionar host de build; instalar dependencias de \textbf{Digi Embedded Yocto (DEY 4.0)}.
  \item \file{repo}/\file{git clone} de las capas \file{meta-digi}; \file{mkproject.sh}
        para plataforma \textbf{\file{ccmp25}} (ConnectCore MP255).
  \item Lanzar el \textbf{primer} \file{bitbake <dey-image-*>} (tarda \textbf{horas};
        dejar corriendo toda la noche).
\end{itemize}

% ───────────
\subsection{Día 2 — 5G como segundo underlay + primer boot embebido}

\textbf{Objetivo:} WAN de respaldo lista y la imagen custom arrancando en el MP255.

\textbf{Pasos (in-band):}
\begin{enumerate}
  \item Configurar el módem 5G del IX40 con SIM/APN de Carrier A → segundo underlay in-band.
  \item Métricas WAN: \file{ETH1} métrica IPv4 = 1 (Starlink, prioritaria), \file{Modem} = 3.
  \item SureLink \textbf{básico} en ambas interfaces (sin afinar todavía).
  \item Registrar IPv4/IPv6/MTU/NAT del 5G (\textbf{\gate{10}}).
\end{enumerate}

\begin{clave}
\small
\done{} al desconectar Starlink a mano, el tráfico sale por 5G (verificación gruesa,
sin medir aún).
\end{clave}

\textbf{\embebido}
\begin{itemize}
  \item Flashear la imagen del Día 1 en el MP255; \textbf{primer boot}.
  \item Validar arranque, consola, red del SOM. Guardar la imagen base reproducible.
\end{itemize}

% ───────────
\subsection{Día 3 — Fibra de servicio + overlay + eBGP}

\textbf{Objetivo:} el sitio anuncia su red por BGP sobre el overlay al core.

\textbf{Pasos (in-band):}
\begin{enumerate}
  \item Poblar \file{SFP/ETH5} con la óptica \textbf{1\,Gb/s validada (\gate{6})} →
        fabric/carga de servicio. \textbf{Nunca SFP+.}
  \item Definir la red de servicio del sitio: IPv4 RFC1918 \file{/24} + IPv6 documentación
        \file{/64} (\file{2001:db8::/32}).
  \item Levantar el \textbf{overlay IKEv2/IPsec route-based con NAT-T} hacia el endpoint
        central público (NAT-T obligatorio por CGNAT).
  \item Configurar \textbf{eBGP} sitio (AS65001) ↔ core (AS65000) bajo
        \file{network route service bgp}. Capturar el nombre real del campo de \textbf{ASN remoto}
        del vecino inspeccionando el esquema DAL (\file{?} en
        \file{network route service bgp neighbour} o la API local
        \file{https://\{ix40\}/cgi-bin/config.cgi}) — \textbf{\gate{4}}, no adivinar.
  \item Aplicar políticas: prefix-lists in/out, max-prefix, prohibir anunciar default.
\end{enumerate}

\begin{clave}
\small
\done{} el collector AS65000 recibe \textbf{exactamente} los prefijos IPv4/IPv6 esperados del
sitio; los filtros negativos pasan (no se filtra lo prohibido); MTU documentada;
cero route leaks.
\end{clave}

\textbf{\embebido}
\begin{itemize}
  \item Crear la \textbf{receta custom} en una capa propia (\file{meta-<lab>}) que empaquete el
        \textbf{agente} (Python) en la imagen.
  \item Rebuild incremental (mucho más rápido que el primero).
\end{itemize}

% ───────────
\subsection{Día 4 — DRM-as-code + alarmas + Push Monitor → SIEM}

\textbf{Objetivo:} el sitio se gestiona como código y emite eventos al SIEM.

\textbf{Pasos:}
\begin{enumerate}
  \item \textbf{Enrolar} el IX40 por API:
        \file{POST /ws/v1/devices/inventory?bundle\_device=true} con \file{id} +
        \file{install\_code} (de la etiqueta, nunca en el repo).
  \item Importar la config del IX40 de referencia, leer el árbol real y crear la
        \textbf{golden config} con Configuration Manager:
        \file{POST /ws/v1/configs/inventory}
        (campos \file{name}, \file{groups}, \file{type}, \file{vendor\_id}, \file{enabled},
        \file{alert}, \file{remediate}, \file{scan\_frequency},
        \file{maintenance\_window\_handling}…). \textbf{Obtener \file{type}/\file{vendor\_id} del
        IX40 real (\gate{12}); no copiar el ejemplo TX64.}
  \item Settings comunes vs por sitio (\file{/settings} y
        \file{/settings/device?device\_id=…}); recordar que un \file{PUT} sustituye el árbol
        enviado → declarar solo nodos administrados (\file{@managed}, \file{@scope},
        \file{\#value}).
  \item Validar el body de \file{scan\_now} contra \file{GET /ws/v1/configs} o el API Explorer
        del tenant (\textbf{\gate{3}}); iniciar un scan y consultar
        \file{GET /ws/v1/configs/compliance}.
  \item Alarmas: \file{POST /ws/v1/alerts/inventory} (Device Offline con
        \file{fire.parameters.reconnectWindowDuration}) + noncompliance.
  \item \textbf{Push Monitor:} \file{POST /ws/v1/monitors/inventory}
        (\file{type: http}, \file{topics: [alert\_status]}, \file{url} al SIEM,
        \file{persistent: true}).
\end{enumerate}

\begin{clave}
\small
\done{} flujo completo observado con contratos validados: enrolamiento, golden config aplicada,
scan/compliance, alarma disparada y evento recibido por el SIEM vía Push Monitor.
\end{clave}

\textbf{\embebido}
\begin{itemize}
  \item El agente recolecta \textbf{red local + ambientales} y envía al SIEM
        (formato estructurado, mismo reloj NTP, ID por evento).
\end{itemize}

% ───────────
\subsection{Día 5 — SureLink multicapa + primera campaña de failover nativo}

\textbf{Objetivo:} failover afinado y \textbf{medido} (no ``ping ok'').

\textbf{Pasos:}
\begin{enumerate}
  \item SureLink con pruebas \textbf{multicapa} e independientes: \file{interface\_up},
        ping de gateway, ping externo, \file{dns}, \file{https} y equivalentes IPv6.
        Múltiples destinos (evitar punto único de fallo).
  \item Acciones \textbf{escalonadas}: subir métrica → reiniciar interfaz; reset de
        módem/cambio de SIM solo tras fallos persistentes.
  \item \textbf{Histéresis + hold-down} contra flapping; \textbf{failback} a Starlink solo
        tras ventana estable.
  \item Montar el testigo de medición: ping continuo + una videollamada/stream como carga real.
  \item Ejecutar la \textbf{campaña representativa} (≈10 reps): desconexión física de Starlink
        y blackhole aguas arriba (Ethernet up). Correlacionar cada inyección con un ID único.
\end{enumerate}

\begin{clave}
\small
\done{} para desconexión y blackhole, se tiene tiempo de detección, tiempo de conmutación y
pérdida por repetición, con p50/p95; trazas correlacionadas con la alerta DRM.
\end{clave}

\textbf{\embebido}
\begin{itemize}
  \item Implementar \textbf{buffer offline + reenvío} del agente: al cortar la red, sigue
        recolectando; al reconectar, reenvía sin perder datos.
\end{itemize}

\begin{nota}
\small
\stretchgoal{Semana 1:} detección ligera de anomalías en el agente (umbral/estadística simple sobre
las series recolectadas).
\end{nota}

% ─────────────────────────────────────────────────────────────────
\section{Semana 2 — Resiliencia avanzada, embebido a producción, clonado, pruebas, demo}

% ───────────
\subsection{Día 6 — WAN Bonding A/B + endurecimiento del agente embebido}

\textbf{Objetivo:} comparar los dos mecanismos de resiliencia y endurecer el agente embebido
(secure boot se difiere a fase posterior).

\textbf{Pasos (in-band):}
\begin{enumerate}
  \item Aprovisionar el \textbf{servidor externo de WAN Bonding} (Tier 3, 1\,Gb/s) con
        credenciales de túnel.
  \item En el IX40: \file{network sdwan wan\_bonding}, meter \file{ETH1} y \file{Modem} en el
        túnel, cifrado on. Estado con \file{show wan-bonding}.
  \item Reproducir \textbf{el mismo impairment trace} del Día 5 y comparar native vs bonding:
        continuidad de sesión, throughput, RTT, jitter, pérdida.
\end{enumerate}

\begin{clave}
\small
\done{} informe comparativo native vs bonding con el mismo trace; la sesión de la app de prueba
\textbf{no} se resetea bajo bonding (objetivo pérdida ≤ 1\,\% en handoff).
\end{clave}

\textbf{\embebido}
\begin{itemize}
  \item \textbf{Endurecer el agente + validar resiliencia:} afinar la detección de anomalías y
        repetir cortes de red bajo carga para confirmar buffer/replay sin pérdida.
  \item \textit{TrustFence secure boot se movió a fase posterior (F2-C) por ser irreversible
        (\file{trustfence close}) y de tiempo impredecible.}
\end{itemize}

% ───────────
\subsection{Día 7 — Telemetría XBee + Digi Container + arranque de OTA}

\textbf{Objetivo:} cerrar la telemetría física y el edge compute del router.

\textbf{Pasos:}
\begin{enumerate}
  \item Emparejar el sensor \textbf{XBee 3 Zigbee} con el \textbf{Hive}; publicar la lectura
        (temperatura/puerta/entrada digital) como \textbf{DataPoints} en DRM.
  \item Desplegar un \textbf{Digi Container} ligero (collector de métricas) al IX40 desde DRM,
        con \textbf{límites de recursos}.
  \item Prueba de carga del container: demostrar que \textbf{no degrada} el plano de forwarding.
\end{enumerate}

\begin{clave}
\small
\done{} la telemetría llega por ambos underlays; uso de recursos del container medido y sin
impacto en el forwarding.
\end{clave}

\textbf{\embebido}
\begin{itemize}
  \item Configurar \textbf{OTA} vía ConnectCore Cloud Services: \file{/etc/cccs.conf},
        construir un paquete \file{.swu}, subirlo al \textbf{firmware repository de Remote
        Manager} (p.\,ej.\ con \file{python-devicecloud}) y disparar una actualización de
        prueba. \textit{(Stretch del sprint; completo en Anillo 2.)}
\end{itemize}

% ───────────
\subsection{Día 8 — Clonado de sitio con automatización}

\textbf{Objetivo:} demostrar que un sitio nuevo se levanta con un comando.

\textbf{Pasos:}
\begin{enumerate}
  \item Modelar la config como \textbf{datos}: plantillas Jinja2 + variables por sitio
        (ASN, prefijos, tunnel peer, nombres, carrier).
  \item Envolver todo en Ansible/Python + \file{make deploy}: enrolamiento API → golden config
        por grupo → settings por sitio → verificación.
  \item Levantar el \textbf{Sitio B}:
        \begin{itemize}
          \item Si llegó el hardware → sitio \textbf{físico} clonado.
          \item Si no → clonar contra un dispositivo de evaluación/simulado y
                \textbf{documentar} que la limitación es de hardware, no de método.
        \end{itemize}
  \item Probar \textbf{idempotencia}: correr \file{make deploy} dos veces no rompe nada.
\end{enumerate}

\begin{clave}
\small
\done{} un segundo sitio gestionado surge del comando; el registro de auditoría
(vía RADIUS/TACACS+ → SIEM) muestra cada acceso trazado; idempotencia demostrada.
\end{clave}

\begin{nota}
\small
\stretchgoal{GitOps} — un cambio en Git dispara la reconciliación de la config de la ``flota''.
\end{nota}

% ───────────
\subsection{Día 9 — Campaña de caos + seguridad + congelar datos}

\textbf{Objetivo:} ejercer un subconjunto de la matriz de fallos y dejar la evidencia limpia.

\textbf{Pasos:}
\begin{enumerate}
  \item Ejecutar los fallos del sprint (ver \file{06-pruebas-y-demo.md} §Matriz):
        desconexión Starlink, blackhole, DNS-only, drift de config, receiver de webhook offline.
  \item Correlacionar \textbf{alarma ↔ evento Push Monitor ↔ ID de fallo}; probar
        \textbf{replay} desde \file{GET /ws/v1/monitors/history/\{monitorId\}} tras dejar el
        receiver caído.
  \item Drift + \textbf{remediación} por Configuration Manager dentro de una maintenance window.
  \item Revisión de seguridad: usuarios de aplicación separados, API keys con vencimiento,
        MFA/RBAC, cero defaults, \textbf{cero secretos} en artefactos; sanitizar evidencia.
\end{enumerate}

\begin{clave}
\small
\done{} todos los casos del sprint tienen resultado, evidencia y dueño; replay comprobado;
drift remediado; datasets congelados con versión de firmware/config.
\end{clave}

% ───────────
\subsection{Día 10 — Empaquetado y demo}

\textbf{Objetivo:} dejar el capstone entregable.

\textbf{Pasos:}
\begin{enumerate}
  \item Diagrama \textbf{lógico} y \textbf{as-built}; BOM validada; documento de intención de
        configuración por firmware.
  \item Contrato de API DRM validado contra el tenant real (y Lighthouse en fase 2).
  \item Runbook repetible (onboarding, failover, OOB, rollback, recuperación).
  \item Grabar la \textbf{demo de 12 min} + versión de \textbf{90 s}; ensayar
        (ver guion en \file{06-pruebas-y-demo.md} §Demo).
  \item \textbf{One-pager ejecutivo} para el manager + \textbf{README sanitizado} para
        portafolio.
\end{enumerate}

\begin{clave}
\small
\done{} demo de 12 min ensayada; paquete entregado al manager; versión pública sanitizada lista;
backlog de gaps (G1–G12) con dueño recomendado.
\end{clave}

% ─────────────────────────────────────────────────────────────────
\section{Fase 2 — Opengear + expansión (cuando llegue el hardware)}

\begin{nota}
\small
Estos bloques \textbf{no} están en el camino crítico del sprint; se ejecutan cuando el hardware
Opengear y las unidades 2-3 llegan. Mapea a las semanas 5–10 de plan2.
\end{nota}

\subsection{F2-A · OOB / Lighthouse}

\begin{enumerate}
  \item Cerrar \textbf{\gate{1}} (confirmar SKU \file{OM1304-4E-C-LSP}, revisión, región,
        firmware) — o activar fallback \file{CM8004-C-LSP} + switch externo.
  \item Enrolar el OM por \textbf{LSP} en Lighthouse Enhance
        (\file{POST /nodes} con \file{enrollment};
        \file{Authorization: Token \{session\}}).
        Probar \textbf{factory reset/boot} para demostrar zero-touch
        (\file{systemctl \{start,enable\} lsp}).
  \item Consola serial del IX40 en modo \textbf{Login} (adaptador \file{319015}), acceso
        restringido por RBAC.
  \item Validar \file{runtime\_status.connection\_status} y \textbf{acceso a consola durante
        pérdida total del plano in-band}.
  \item Secure Provisioning downstream: \textbf{opcional}, hasta confirmar licencia en
        Core/Enhance (\gate{2}/\gate{11}); observar progreso en syslog
        \file{[NetOps-DOP device="MAC"]}.
\end{enumerate}

\subsection{F2-B · Tres sitios + rotación + 30 reps}

\begin{enumerate}
  \item Clonar Sitios B y C físicos con la automatización del Día 8.
  \item Rotar perfiles \textbf{native/bonding/candidate} entre las tres unidades.
  \item Campaña de \textbf{30 repeticiones} por fallo con p50/p95/p99 y datos brutos.
\end{enumerate}

\subsection{F2-C · Embebido a producción}

\begin{enumerate}
  \item \textbf{TrustFence secure boot:} generar claves, \file{SRK\_efuses.bin},
        firmar U-Boot + kernel, \file{trustfence close}
        (\textbf{irreversible}; practicar en unidad sacrificable),
        cifrado de rootfs (\file{CONFIG\_FS\_ENCRYPTION}, datos en \file{/mnt/data/private}).
  \item \textbf{OTA completo} del agente (\file{.swu} por CCCS) sin ladrillar el dispositivo.
  \item Carril \textbf{OEM} con el MP255 (TSN/NPU) en Site C.
  \item Lazo de \textbf{auto-remediación}: árbol de decisión según tipo de fallo,
        verificación post-remediación y rollback.
\end{enumerate}

\begin{aviso}
\small
\textbf{\file{trustfence close} es irreversible.} Ejecutar únicamente con la imagen ya estable
y en una unidad sacrificable de prueba antes de aplicar a producción.
\end{aviso}

% ─────────────────────────────────────────────────────────────────
\section{Notas de secuenciación}

\begin{itemize}
  \item \textbf{Yocto arranca el Día 1} porque el primer build tarda horas; si esperara a la
        semana 2, bloquearía el track central.
  \item \textbf{DRM-as-code (Día 4) antes del clonado (Día 8):} no puedes clonar lo que aún no
        modelaste como datos.
  \item \textbf{SureLink medido (Día 5) antes de Bonding (Día 6):} necesitas el baseline nativo
        para que la comparación A/B signifique algo.
  \item \textbf{\file{trustfence close} es irreversible:} por eso \textbf{no} entra en el sprint;
        se ejecuta en fase posterior (F2-C) con la imagen ya estable.
  \item \textbf{Opengear fuera del camino crítico:} su retraso no puede frenar el sprint;
        entra completo en fase 2.
\end{itemize}

\end{document}
